\documentclass{article}
\usepackage{pathfinder-note}

\title{From Distributed CVaR Optimization to a Costed Risk Certificate}

\begin{document}
\maketitle

\section*{The pair}

Q surveys learning-augmented algorithms and argues that prediction costs, benchmark relations, and downstream interfaces need explicit accounting before component guarantees can be transferred to a system. P proves finite-time expected CVaR suboptimality and exact decision-domain feasibility for a distributed algorithm using sampled local losses, but its objective is the average of local CVaRs. \ledger{3, 11, 13}

\section*{The connection pursued}

The question was whether P's output can support a verifiable, costed guarantee for the aggregate loss incurred when its decision is deployed. Write \(X=\hat x_T\) for P's weighted ergodic output, \(C_i(x)\) for agent \(i\)'s CVaR, and
\[
\mathcal C(x)=\frac{1}{m}\sum_{i=1}^m C_i(x).
\]
Let \(B_T\) denote a deterministic upper bound on P's finite-time result:
\[
\mathbb E[\mathcal C(X)-\mathcal C^*]\leq B_T.
\]
P's projection makes \(X\in\mathcal X\) pathwise. The expectation bound alone neither certifies the risk of a realized \(X\) nor establishes feasibility for a separate downstream risk constraint. \ledger{3, 13}

\section*{What was established}

Assume a common tail level \(\alpha\), and suppose deployment incurs the same agents' losses jointly:
\[
S(x,\xi)=\sum_{i=1}^m J_i(x,\xi_i),
\qquad
A(x)=\operatorname{CVaR}_{\alpha}(S(x,\xi)).
\]
CVaR subadditivity gives \(A(x)\leq m\mathcal C(x)\). With \(A^*=\min_{x\in\mathcal X}A(x)\) and the comparator gap
\[
D=m\mathcal C^*-A^*\geq0,
\]
P's theorem consequently implies
\[
\mathbb E[A(X)-A^*]\leq mB_T+D.
\]
This is the benchmark relation required by Q's composition discussion. It does not bound \(D\) from P's local optimization error. \ledger{6, 22}

The gap can be actual decision loss. For \(m=2\), \(\alpha=1/2\), \(x\in[-1,1]\), and a shared fair Bernoulli variable \(Z\), take
\[
J_1(x,Z)=UZ,\qquad
J_2(x,Z)=U\bigl((1-x)Z+x(1-Z)\bigr).
\]
The sum of local CVaRs is minimized at \(x_P=1/2\), whereas \(A(x)=U(2-x)\) is minimized at \(x_A=1\). Thus \(A(x_P)-A^*=D=U/2\), even with exact local optimization. \ledger{23, 27}

A selective certificate is possible under additional assumptions. Freeze \(X\), then draw independent validation samples at that decision. If agent \(i\) supplies \(n_i\) fresh samples bounded in absolute value by \(U_i\), P's CVaR stability lemma and the DKW inequality give simultaneous errors at most
\[
\epsilon_i=
\frac{2U_i}{\alpha}
\sqrt{\frac{\log(2m/\zeta)}{2n_i}}
\]
with probability at least \(1-\zeta\). On that validation event,
\[
A(X)\leq
R:=\sum_{i=1}^m\bigl(\widehat C_i(X)+\epsilon_i\bigr).
\]
Hence accepting \(X\) only when \(R\leq b\) has probability at most \(\zeta\) of accepting an action whose aggregate CVaR exceeds \(b\). For a fresh deployment draw, \(A(X)\leq R\) also implies
\(\Pr_\xi\{S(X,\xi)>R\}\leq\alpha\).
The confidence probability over validation data and the tail probability over future loss are distinct. \ledger{4, 6, 13}

With a known safe fallback \(x_0\), a known margin
\(m\mathcal C^*\leq b-h\), and validation precision
\(2\sum_i\epsilon_i\leq h/2\), the peers derived
\[
\Pr\{\text{reject }X\}
\leq \zeta+\frac{2mB_T}{h}.
\]
Writing \(\Delta=\mathcal C(x_0)-\mathcal C^*\), the expected deployed local-objective gap is at most
\[
B_T+\Delta\left(\zeta+\frac{2mB_T}{h}\right).
\]
The margin and safe fallback are extra premises, not consequences of P. \ledger{8, 14}

\section*{Objections and limits}

The local certificate can be arbitrarily conservative in the number of agents: mutually exclusive losses can make \(\sum_i C_i=mU\) while their aggregate loss is always \(U\). Identical local marginals can also have different aggregate CVaRs under different couplings. A marginal-only oracle cannot identify that dependence, even at P's exact optimizer; synchronized joint observations or a justified dependence model are needed. \ledger{18, 21, 26}

A fresh paired audit of the actual aggregate loss removes this certificate slack at the chosen \(X\). With \(n\) independent paired loss vectors, \(U_\Sigma=\sum_iU_i\), and empirical aggregate CVaR \(\widehat A(X)\), the DKW argument gives
\[
A(X)\leq R_J:=\widehat A(X)+
\frac{2U_\Sigma}{\alpha}
\sqrt{\frac{\log(2/\zeta)}{2n}}
\]
with confidence \(1-\zeta\). It still does not identify \(A^*\) or remove \(D\) from the decision-quality bound. \ledger{22, 28}

Finite samples do not provide exact distribution-free risk feasibility for every accepted action. The ledger constructs bounded, convex losses with a safe fallback for which a safe law and an unsafe law both produce the same finite all-zero transcript with positive probability. A nontrivial acceptance rule can therefore make a false acceptance; the confidence term cannot simply be dropped. \ledger{9, 10}

Q prices predictor invocations, while P queries sampled loss values and communicates decisions. Charging \(N\) such evaluations as \(\kappa N\) is a new accounting convention. If a query physically deploys a perturbed action, its incurred loss must also be charged; producing \(X\), pairing observations, and communicating audit results can carry further costs. No competitive ratio for a separate downstream policy follows without its own sensitivity and benchmark relation. \ledger{11, 26, 28}

\section*{Outside sources and remaining question}

The peers checked Liang and Luo's work on CVaR confidence bounds at
\url{https://proceedings.mlr.press/v202/liang23c.html}
and Cheung's work on dependence in risk aggregation at
\url{https://doi.org/10.1016/j.insmatheco.2010.06.001}.
The holdout bound and dependence obstruction are therefore a conditional bridge between Q and P, not a demonstrated novelty claim. \ledger{15, 19}

What remains open is a useful cost frontier for local training, joint observation, validation, and fallback under a specified deployment interface. The smallest next step is to fix one application with a priced paired-sample oracle and a safe baseline, then establish or refute a usable bound on \(D\) and compare the resulting total cost with that baseline. The present material does not establish such an improvement. \ledger{24, 28}

\end{document}